\documentclass{subfiles}
\begin{document}
    The core idea of the method follows that used by banks.
        In more details: to each operation we assign a cost,
        which represents its \emph{amortized cost}.
        When the amortized cost of an operation exeeds its real cost, 
        we are left with some credit that we can use to help pay other operations.

    \begin{example*}
        Consider we need an algorithm to increment a \(n\) bits counter by 
            an amount \(m\). From a worst case analysis perspective,
            the algorithm needs a quadratic time, as the \(n\)-th bit may need to be flipped \(m\) times.
            It comes with no surprise that such a cost is well above the real one,
            as we are about to see.

        Let us proceed this way: for each bit flip, assume that we are asked to pay 
            a coin; then we proceed by paying two coins when flipping a bit to 1.
            The cost of the whole algorithm is easily seen to be given by the number of 
            total bit flipped. But the following can be observed: the first increment,
            leaves us with a coin of credit, the second one does the same,
            etc. Thus, each operation requires a constant amount;
            that is, the amortized cost of the whole algorithm is \(\ofO{n}\).            
    \end{example*}

    \begin{remark*}
        Other methods to perform amortized analysis exist, 
            and their use varies with the algorithm we need to analyze.
            Though, all follow the same core idea: we overpay for some operations
            and use the credit thus produced to pay for others.
    \end{remark*}
\end{document}
